\subsection*{6. Constructibilidad de $\R\underline{\Hom}$}

La equivalencia de las condiciones a) de (I 3.3.1) sigue siendo válida cuando se sustituye el anillo $(\Z/n\Z)_X$ por un anillo conmutativo, localmente constante y cuyas fibras geométricas sean anillos noetherianos: no hay que modificar nada en la demostración.

Las condiciones anteriores se verifican cuando se dispone de la resolución de singularidades y de la pureza en sentido fuerte (4.4), en particular si $X$ es de dimensión $\leq1$, si $X$ es excelente de característica cero, o si es localmente de tipo finito sobre un cuerpo y de dimensión $\leq2$.

\subsection*{7. Existencia de complejos dualizantes}

En este número, $X$ designa un esquema provisto de un haz de anillos $A_X$ conmutativo, localmente constante y cuyas fibras geométricas son anillos noetherianos.

\textbf{Lema 7.1.} Sean $F\in\ob\D^{-}_{\mathrm{coh}}(X)$ y $G\in\ob\D^{+}_{\mathrm{coh}}(X)$; entonces $\R\underline{\Hom}(F,G)\in\ob\D^{+}_{\mathrm{coh}}(X)$ (notaciones de (1.2)).

\emph{Demostración.} Es estándar. Mediante una sucesión espectral, se reduce al caso en que $F$ y $G$ están concentrados en grado $0$, y después se procede como en (EGA $0_{\mathrm{III}}$ 12.3.3).

\textbf{Lema 7.2.} Sea $g:X'\to X$ un morfismo de esquemas, y dotemos a $X'$ del haz de anillos $A_{X'}=g^*(A_X)$. Entonces, para $F\in\ob\D^{-}_{\mathrm{coh}}(X)$ y $G\in\ob\D^{+}(X)$, la flecha canónica
\[
        g^*\R\underline{\Hom}(F,G)\longrightarrow \R\underline{\Hom}(g^*F,g^*G)
\]
es un isomorfismo.

\emph{Demostración.} También es estándar: por devissage, se reduce al caso $F=A_X$ (cf.~EGA $0_{\mathrm{III}}$ 12.3.5).

\textbf{Observación 7.2.1.} Los lemas anteriores son casos particulares de enunciados más generales que figuran en (SGA 6 I).

\textbf{Lema 7.3.} Supongamos que $X$ es cuasicompacto, y sea $K\in\ob\D^{+}_{\mathrm{coh}}(X)$ tal que, para todo punto geométrico $\bar{x}$ de $X$, $K_{\bar{x}}$ sea un complejo dualizante en el sentido de [H] V 2.1. Entonces, para todo $F\in\ob\D^{b}_{\mathrm{coh}}(X)$, se tiene $\R\underline{\Hom}(F,K)\in\ob\D^{b}_{\mathrm{coh}}(X)$, y $F$ es reflexivo con respecto a $K$.

\emph{Demostración.} Naturalmente, basta efectuar la verificación cuando $F$ está concentrado en grado $0$. Por (7.1), se tiene $\R\underline{\Hom}(F,K)\in\ob\D^{+}_{\mathrm{coh}}(X)$. Demostremos que $\R\underline{\Hom}(F,K)\in\ob\D^{b}_{\mathrm{coh}}(X)$. Sea $\bar{x}$ un punto geométrico de $X$. Por (7.2), se tiene
\[
        \R\underline{\Hom}(F,K)_{\bar{x}}\simeq \R\underline{\Hom}(F_{\bar{x}},K_{\bar{x}}).
\]
Si ponemos $\dim.\mathrm{inj}(K_{\bar{x}})=N(\bar{x})$, se tiene, por tanto,
\[
        \underline{\Ext}^i(F,K)_{\bar{x}}=0\quad\hbox{para }i>N(\bar{x}).
\]
Como los $\underline{\Ext}^{i}(F,K)$ son localmente constantes, existe un entorno abierto $U$ de $x$ tal que $\underline{\Ext}^{i}(F,K)|U=0$ para $i>N(\bar{x})$, y se concluye gracias a la hipótesis de cuasicompacidad sobre $X$. Resta comprobar que la flecha canónica
\[
        F\longrightarrow \R\underline{\Hom}(\R\underline{\Hom}(F,K),K)
\]
es un isomorfismo. Para ello basta situarse en un punto geométrico $\bar{x}$ de $X$. En el diagrama conmutativo
\[
\begin{tikzcd}
F_{\bar{x}} \arrow[r] \arrow[d,"\mathrm{Id}"'] & \R\underline{\Hom}(\R\underline{\Hom}(F,K),K)_{\bar{x}} \arrow[d]\\
F_{\bar{x}} \arrow[r] & \R\underline{\Hom}(\R\underline{\Hom}(F_{\bar{x}},K_{\bar{x}}),K_{\bar{x}}),
\end{tikzcd}
\]
la flecha vertical de la derecha es un isomorfismo por (7.2), y la flecha horizontal inferior es un isomorfismo porque $K_{\bar{x}}$ es dualizante; por tanto, la flecha horizontal superior es un isomorfismo, c.q.d.

\textbf{Lema 7.4.} Supongamos que $X$ es un esquema noetheriano regular, que $A_X$ está anulado por un entero $n>0$ primo con las características residuales de $X$, y que el teorema de pureza en sentido fuerte (4.4) es válido sobre $X$.

Sea $D$ un divisor de cruzamientos normales sobre $X$ (I 3.1.5 a)); denotemos por $j:Y=V(D)\to X$ e $i:X-Y=U\to X$ las inclusiones canónicas. Sea $K\in\ob\D^{+}_{\mathrm{coh}}(X)$ tal que, para todo punto geométrico $\bar{x}$ de $X$, $K_{\bar{x}}$ sea un complejo dualizante para el anillo $A_{\bar{x}}$ (en el sentido de [H] V 2.1). Entonces, para todo $E\in\ob\D^{b}_{\mathrm{coh}}(X)$, se tiene $\R\underline{\Hom}(i_!i^*(E),K)\in\ob\D^{b}_{c}(X)$, y el par $(i_!i^*(E),K)$ es bidualizante.

\emph{Demostración.} Podemos suponer que $E$ está concentrado en grado $0$, es decir, que es un $A_X$-módulo localmente constante constructible. Por otra parte, como la cuestión es local sobre $X$ para la topología étale, podemos suponer que
\begin{enumerate}[label=(\alph*)]
\item $A_X$ es constante de valor $A$, y $E$ es constante de valor denotado también por $E$, por abuso de lenguaje; es decir, $E=E_X$;
\item $D$ es de cruzamientos normales estrictos, es decir, $D=\sum_{i\in I}\mathrm{div}(s_i)$, donde $s_i\in\Gamma(X,\mathcal{O}_X)$.
\end{enumerate}
Razonemos por inducción sobre $p=\operatorname{card}(I)$. Para $p=0$, las afirmaciones son consecuencia de (7.3). Supongamos que $I=\{1,\ldots,p\}$, pongamos $Y_i=V(s_i)$, $Y'=V(\sum_{1\leq i\leq p-1}\mathrm{div}(s_i))$, $U'=X-Y'$; se obtiene así un diagrama cartesiano de inmersiones
\[
\begin{tikzcd}
\varnothing \arrow[r] \arrow[d] & U'\cap Y_p \arrow[r,"i'_p"] \arrow[d,"j'_p"'] & Y_p \arrow[d,"j_p"]\\
U \arrow[r] & U' \arrow[r,"i'"'] & X.
\end{tikzcd}
\]
Consideremos la sucesión exacta
\[
        0\longrightarrow E_{U,X}\longrightarrow E_{U',X}\longrightarrow E_{U'\cap Y_p,X}\longrightarrow0.
\]
Por hipótesis de inducción, el teorema es válido para el par $(E_{U',X},K)$. Como, para $K$ fijo, el conjunto de los $F$ tales que $\R\underline{\Hom}(F,K)\in\ob\D_c^b(X)$ y que $(F,K)$ sea bidualizante forma una subcategoría triangulada de $\D_c(X)$, quedamos reducidos a demostrar el teorema para el par $(E_{U'\cap Y_p,X},K)$. Pero se tiene
\[
        E_{U'\cap Y_p,X}=j_{p*}i'_{p!}(E_{U'\cap Y_p}),
\]
y, por (I 1.13), equivale a demostrar el teorema para el par $(i'_{p!}(E_{U'\cap Y_p}),\R j_p^!(K))$ en $Y_p$. Pero, por el teorema de pureza, se tiene (4.3)
\[
        \R j_p^!(K)\simeq j_p^*(K)\otimes^L_{\Z/n\Z}(\mu_n)_{Y_p}^{\otimes -1}[-2].
\]
Por tanto, $\R j_p^!(K)\in\ob\D^{+}_{\mathrm{coh}}(Y_p)$ y, para todo punto geométrico $\bar{y}$ de $Y_p$, $(\R j_p^!(K))_{\bar{y}}$ es dualizante; se concluye entonces por la hipótesis de inducción.

\textbf{Teorema 7.5.} Hagamos las hipótesis siguientes: $X$ es un esquema noetheriano regular; $A_X$ está anulado por un entero $n>0$ primo con las características residuales de $X$; $X$ es fuertemente desingularizable (I 3.1.5.), satisface las condiciones (reg) (I 3.1.1) y $(\mathrm{cdloc})_n$ (I 3.1.3); el teorema de pureza en sentido fuerte (4.4) es válido sobre todo esquema liso sobre $X$ (condiciones que se satisfacen si $X$ es excelente de característica cero, o si es de dimensión $\leq2$ y liso sobre un cuerpo perfecto).

Entonces, si $K\in\ob\D^{+}_{\mathrm{coh}}(X)$ es tal que $K_{\bar{x}}$ sea dualizante (en el sentido de [H] V 2.1) para todo punto geométrico $\bar{x}$ de $X$, $K$ es dualizante.

\emph{Demostración.} Se trata de demostrar que $K$ satisface las condiciones (i) a (iii) de (I 1.7). Para todo punto geométrico $\bar{x}$ de $X$, $K_{\bar{x}}$ tiene dimensión inyectiva finita y, como $K\in\ob\D^{+}_{\mathrm{coh}}(X)$, la función $x\mapsto \dim.\mathrm{inj}(K_{\bar{x}})$ es localmente constante y, por tanto, acotada, pues $X$ es noetheriano. Se deduce entonces de (5.1) que $K$ tiene dimensión cuasi-inyectiva finita.

Por lo visto en el n\textsuperscript{o}~6, también se verifica la condición (ii) de loc.~cit.

Resta demostrar la condición (iii), es decir, que todo $F\in\ob\D_c^b(X)$ es reflexivo con respecto a $K$. Por devissage, podemos suponer que $F$ está concentrado en grado $0$. Por otra parte, como la cuestión es local sobre $X$ para la topología étale, podemos suponer que $A_X$ es constante de valor $A$. El devissage (1.3) nos reduce entonces a $F=f_*i_!(E_U)$, donde $f:Y\to X$ es un morfismo finito, con $Y$ íntegro y su punto genérico separable sobre su imagen, $i:U\to Y$ es un abierto regular no vacío, y $E_U$ es un $A_U$-módulo constante de valor $E$, un $A$-módulo de tipo finito. Como $X$ es fuertemente desingularizable, existe (I 3.1.5) un morfismo proyectivo y birracional $h:Y'\to Y$ tal que $Y'$ sea regular, que $h$ induzca un isomorfismo $h^{-1}(U)\to U$, y que $h^{-1}(U)$ sea el complementario de un divisor de cruzamientos normales en $Y'$. Identificando $h^{-1}(U)$ con $U$ mediante $h$, se tiene el diagrama conmutativo
\[
\begin{tikzcd}
& Y' \arrow[d,"h"] \arrow[dr,"g"] & \\
U \arrow[ur,"i'"] \arrow[r,"i"'] & Y \arrow[r,"f"'] & X.
\end{tikzcd}
\]
Se tiene, por consiguiente,
\[
        f_*i_!(E_U)\simeq f_*\R h_* i'_!(E_U)\simeq \R g_* i'_!(E_U)\simeq \R g_* i'_! i'^*(E_{Y'}).
\]
En virtud de (I 1.13), basta demostrar la bidualidad para el par $(i'_!i'^*(E_{Y'}),\R g^!(K))$. Pero, embebiendo localmente $Y'$ en un esquema $X'$ liso sobre $X$ y aplicando la pureza en sentido fuerte sobre $X'$ (4.4), se ve que $\R g^!(K)$ es localmente isomorfo (salvo traslación de los grados) a $g^*(K)$; por tanto, $\R g^!(K)\in\ob\D^{+}_{\mathrm{coh}}(Y')$ y $\R g^!(K)_{\bar{x}}$ es dualizante para todo punto geométrico $\bar{x}$ de $Y'$. Basta entonces aplicar (7.4).

La demostración anterior prueba, además, que $\R\underline{\Hom}(F,K)\in\ob\D_c^b(X)$ (supuesto el teorema de finitud para un morfismo propio (SGAA XIV)), lo que tranquilizará al lector que no hubiese confiado en la perorata del n\textsuperscript{o}~6.

Esto concluye la demostración.

\begin{center}
$-:\!-:\!-$
\end{center}

\section*{Bibliografía}

\begin{enumerate}[label={[\arabic*]}]
\item Verdier, J.-L., \emph{Dimension des espaces localement compacts}, Note aux C.R. Ac. Sc. Paris, 261, p.~5293--5296 (1965). Véase también: Verdier, J.-L., \emph{Dualité dans la cohomologie des espaces localement compacts}, Séminaire Bourbaki, n\textsuperscript{o}~300, 1965--66.
\end{enumerate}

SGA A = SGA 4.

\begin{center}
$-:\!-:\!-$
\end{center}
